\author{Brett Reynolds}
\title{Mga Tanawin ng Wika: Language Landscapes}
\subtitle{Gabay sa Kayarian ng Ingles para sa mga Guro ng Ingles [Edisyong Tagalog]}
% \ISBNdigital{}
% \ISBNhardcover{}
% \BookDOI{}
\typesetter{Brett Reynolds}
% \proofreader{}
% \lsCoverTitleSizes{51.5pt}{17pt}
\renewcommand{\lsSeries}{tbls}
%\newcommand{\lsSeriesTitle}{Textbooks in Language Sciences}
%\newcommand{\lsSeriesColor}{lsYellow}
\renewcommand{\lsSeriesNumber}{16}
\BackBody{Ang wika ay hindi listahang isinasaulo; isa itong tanawin na dapat siyasatin.

Ang mahusay na naturalista sa larangan ay hindi lamang nagbibigay ng pangalan. Pinagmamasdan niya kung paano kumikilos ang mga bagay, kung saan sila nagsasama-sama, at kung bakit sila nagbabago. Ganito rin ang pagtingin ng aklat na ito sa Ingles. Mas mahalaga rito ang ginagawa ng mga tunay na nagsasalita kaysa ang sinasabi ng mga aklat ng tuntunin tungkol sa dapat nilang gawin. Nakabatay ang balangkas sa \textit{The Cambridge Grammar of the English Language}, at pinasimple ang paliwanag, ngunit hindi ibinaba ang antas.

Umaabante ang mga kabanata mula sa tunog at salita patungo sa parirala, sugnay, at diskurso. May malinaw na halimbawa sa Ingles at maiikling silip sa ibang wika. Sa bawat kabanata, pareho ang gawain: tingnan ang padron, subukin ito sa ebidensiya, pansinin kung saan nasisira ang padron, at itanong kung ano ang itinuturo ng pagkasira. Sa dulo, magkakaroon ka ng magagamit na kaalamang gramatikal para suriin ang mga aklat-aralin, ipaliwanag kung bakit tila kakaiba ang pangungusap ng isang mag-aaral, at ipagtanggol ang sarili mong paliwanag kapag may tanong.

Nasa pagsasanay ka man bilang guro o nagtuturo na sa silid-aralan, pareho ang pakinabang: huminto sa pag-asa sa mga tuntuning walang matibay na batayan, at matutong makita ang sistema ng wika sa sarili mong pagsusuri.}
